% Copyright 2020 by Robert Hildebrand
%This work is licensed under a
%Creative Commons Attribution-ShareAlike 4.0 International License (CC BY-SA 4.0)
%See http://creativecommons.org/licenses/by-sa/4.0/

% Split out of NLP-algorithms.tex on 2026-08-01 (one chapter per file).

\chapter{Computational Issues with NLP}
\todoChapter{ 50\% complete. Goal 80\% completion date: November 20\\
Notes: }
We mention a few computational issues to consider with nonlinear programs.  
\section{Irrational Solutions}
Consider nonlinear problem (this  is even convex)
\begin{align*}
\min \quad & -x\\
\st \quad  & x^2 \leq 2.
\end{align*}
The optimal solution is $x^* = \sqrt{2}$, which cannot be easily represented.  Hence, we would settle for an \textbf{approximate solution} such as $\bar x = 1.41421$, which is feasible since $\bar x^2 \leq 2$, and it is close to optimal.

\section{Discrete Solutions}
Consider nonlinear problem (not convex)
\begin{align*}
\min \quad & -x\\
\st \quad  & x^2 = 2.
\end{align*}
Just as before, the optimal solution is $x^* = \sqrt{2}$, which cannot be easily represented.  Furthermore, the only two feasible solutions are $\sqrt{2}$ and $-\sqrt{2}$.  Thus, there is no chance to write down a feasible rational approximation.

\section{Convex NLP Harder than LP}
Convex NLP is typically polynomially solvable.  It is a generalization of linear programming.  
\begin{general}{Convex Programming}{\polynomial\ \  (typically)}
Given a convex function $f(x)\colon \R^d \to \R$ and convex functions $f_i(x)\colon \R^d \to \R$ for $i=1, \dots, m$,  the \emph{convex programming} problem is
\begin{equation}
\label{eq:convex-programming-comp}
\begin{aligned}
\min \quad & f(x)\\
\st \quad  & f_i(x) \leq 0 & \text{for } i=1, \dots, m\\
& x \in \R^d
\end{aligned}
\end{equation}
\end{general}
\begin{example}
Convex programming is a generalization of linear programming.  This can be seen by letting $f(x) = c^\top x$ and $f_i(x) = A_i x - b_i$.  
\end{example}


\section{NLP is harder than IP}
As seen above, quadratic constraints can be used to create a feasible region with discrete solutions.  For example 
$$
x(1-x) = 0
$$
has exactly two solutions: $x = 0, x=1$.  
Thus, quadratic constraints can be used to model binary constraints.
\begin{general}{Binary Integer programming (BIP) as a NLP}{\nphard}
Given a matrix $A \in \R^{m\times n}$, vector $b \in \R^m$ and vector $c \in \R^n$, the \emph{binary integer programming} problem is
\begin{equation}
\label{eq:BIP}
\begin{aligned}
\max \quad & c^\top x\\
\st \quad  & Ax \leq b\\
& \hcancel[1.5pt]{x \in \{0,1\}^n}\\
& x_i(1-x_i) = 0 & \text{for } i=1, \dots, n
\end{aligned}
\end{equation}
\end{general}


 







\section{Resources}


\begin{resource}
\paragraph{Gradient Free Algorithms - Needler-Mead}
\begin{itemize}
\item \href{https://en.wikipedia.org/wiki/Nelder%E2%80%93Mead_method}{Wikipedia}\\
\item \href{https://youtube/NI3WllrvWoc?t=96}{Youtube}
\end{itemize}
\paragraph{Bisection Method and Newton's Method}
\begin{itemize}
\item See section 4 of the following notes:
\url{http://www.seas.ucla.edu/~vandenbe/133A/133A-notes.pdf}
\end{itemize}
\paragraph{Gradient Descent}
\begin{itemize}
\item \url{https://www.youtube.com/watch?v=tIpKfDc295M}

\item \url{https://www.youtube.com/watch?v=_-02ze7tf08}

\item \url{https://www.youtube.com/watch?v=N_ZRcLheNv0}

\item \url{https://www.youtube.com/watch?v=4RBkIJPG6Yo}
\end{itemize}
\paragraph{Idea of Gradient descent}
\begin{itemize}
\item \url{https://youtu.be/IHZwWFHWa-w?t=323}
\end{itemize}
Vectors:\\
\begin{itemize}
\item \url{https://www.youtube.com/watch?v=fNk_zzaMoSs&list=PLZHQObOWTQDPD3MizzM2xVFitgF8hE_ab&index=2&t=0s}
\end{itemize}
\end{resource}
